\subsection{Coherent coefficients forced by complementary cuts}
\label{sec:peps_dependent_bond_support}

Use the independent incidence spaces and incident representations of
Definitions~\ref{def:peps_dependent_incidence} and
\ref{def:peps_dependent_incident_action}. This section derives a bond-product
expansion from actual cut membership before extracting its coefficients.
The trace duality is \cite[Lemma~4.6]{Schuch2010PEPS}; its use in the
four-block comparison is part of \cite[Theorem~5.5]{Schuch2010PEPS}.
% Source: lemma:noninj:semireg-trace-ug-delta, paper_v3.tex:1015--1029;
% thm:2d:closure, eq:2d:close-in-in, paper_v3.tex:1488--1513.

\begin{definition}[Regrouping physical incidences as bond pairs]%
    \label{def:peps_dependent_bond_product}
    \lean{TNLean.PEPS.DependentBondNetwork.siteBondEquiv,
        TNLean.PEPS.DependentBondNetwork.bondRegrouping,
        TNLean.PEPS.DependentBondNetwork.bondRegrouping_apply,
        TNLean.PEPS.DependentBondNetwork.matrixBondProduct,
        TNLean.PEPS.DependentBondNetwork.bondRegrouping_matrixBondProduct,
        TNLean.PEPS.DependentBondNetwork.representationBondProduct}
    \leanok
    \uses{def:peps_dependent_incidence}
    The bijection $J=SR^{-1}$ sends a site configuration to its head--tail
    bond pairs. It induces the coefficient permutation
    $(\mathcal R\psi)(\beta)=\psi(J^{-1}\beta)$.
    For a matrix family $B$, let
    $b_B(\sigma)=\prod_eB_e((J\sigma)_e^+,(J\sigma)_e^-)$, so
    $(\mathcal Rb_B)(\beta)=\prod_eB_e(\beta_e^+,\beta_e^-)$.
    For group labels $p:\mathcal E\to G$, abbreviate
    $b_p=b_{(U_e(p_e))_e}$.
\end{definition}

\begin{definition}[Product functionals and trace-dual coefficients]%
    \label{def:peps_dependent_bond_coefficient}
    \lean{TNLean.PEPS.DependentBondNetwork.outputBondPairing,
        TNLean.PEPS.DependentBondNetwork.bondCoefficientExtraction}
    \leanok
    \uses{def:peps_dependent_bond_product, thm:peps_torus_projector_extraction}
    Given linear functionals $K_e$ on the bond matrices, define
    \begin{align}
        K(\psi)=\sum_{\beta\in\prod_e(X_e\times X_e)}
        (\mathcal R\psi)(\beta)\prod_eK_e(E_{\beta_e^+,\beta_e^-}).
        \label{eq:peps_dependent_bond_pairing}
    \end{align}
    With the source's trace-dual functionals
    $\ell_{e,g}(B)=\tr[U_e(g^{-1})B\Delta_e]$, write
    $c_p(\psi)=(\bigotimes_e\ell_{e,p_e})(\psi)$ for this product functional.
\end{definition}

\begin{theorem}[Exact extraction of coherent coefficients]%
    \label{thm:peps_dependent_coefficient_extraction}
    \lean{TNLean.PEPS.DependentBondNetwork.outputBondPairing_matrixBondProduct,
        TNLean.PEPS.DependentBondNetwork.bondCoefficientExtraction_bondProduct,
        TNLean.PEPS.DependentBondNetwork.bondCoefficientExtraction_sum,
        TNLean.PEPS.DependentBondNetwork.sum_representationBondProduct_eq_iff}
    \leanok
    \uses{def:peps_dependent_bond_coefficient, def:peps_semiregular}
    For arbitrary $K_e$ and $B_e$, $K(b_B)=\prod_eK_e(B_e)$.
    If every $U_e$ is semi-regular, then
    \begin{align}
        c_p(b_q)                      & =\delta_{p,q},
        \label{eq:peps_dependent_coefficient_delta}    \\
        c_p\left(\sum_q a_qb_q\right) & =a_p.\nonumber
    \end{align}
    Thus two coherent sums in the $b_p$ are equal if and only if their
    complete coefficient functions are equal.
\end{theorem}
\begin{proof}\leanok
    \uses{thm:peps_torus_projector_extraction}
    Distribute (\ref{eq:peps_dependent_bond_pairing}) over the bond pairs.
    The $e$-th sum is $K_e(B_e)$ by expansion in matrix units.
    Trace duality gives $\ell_{e,p_e}(U_e(q_e))=\delta_{p_e,q_e}$;
    their product is $\delta_{p,q}$. Apply linearity to the entire coherent
    sum, then compare each coefficient.
\end{proof}

\begin{definition}[Bond slices and transformed bond matrices]%
    \label{def:peps_dependent_bond_slice}
    \lean{TNLean.PEPS.DependentBondNetwork.bondMatrix,
        TNLean.PEPS.DependentBondNetwork.representation_mem_bondMatrix_range,
        TNLean.PEPS.DependentBondNetwork.physicalBondSlice,
        TNLean.PEPS.DependentBondNetwork.gaugeBondMatrix}
    \leanok
    \uses{def:peps_dependent_bond_product}
    Let $F_e:\C^G\to\C^{X_e\times X_e}$ have columns
    $F_e((a,b),g)=U_e(g)_{a,b}$. Each representation matrix is therefore
    a vector in $\operatorname{im}F_e$.
    The slice at bond $e$ and remaining configuration $\tau$ is
    $S_{e,\tau}\psi(s)=(\mathcal R\psi)(\tau[e\leftarrow s])$.
    Put $B_e(q)=U_e(q_{h(e)})B_eU_e(q_{t(e)}^{-1})$.
\end{definition}

\begin{theorem}[Canonical products, slices, and extracted sums]%
    \label{thm:peps_dependent_canonical_bond_sum}
    \lean{TNLean.PEPS.DependentBondNetwork.dependentProduct_coordinateSlice_eq_smul,
        TNLean.PEPS.DependentBondNetwork.bondRegrouping_network_averagingSite,
        TNLean.PEPS.DependentBondNetwork.network_averagingSite_eq_sum_matrixBondProduct,
        TNLean.PEPS.DependentBondNetwork.bondCoefficientExtraction_network_averagingSite}
    \leanok
    \uses{def:peps_dependent_bond_slice, def:peps_dependent_bond_coefficient}
    A coordinate slice of $\prod_w f_w(\tau_w)$ is
    $(\prod_{w\ne e}f_w(\tau_w))f_e$.
    With $n=|\mathcal V|$, the canonical contraction and its regrouping obey
    \begin{align}
        \mathcal N_{A^0}(B) & =|G|^{-n}\sum_q b_{B(q)},
        \label{eq:peps_dependent_canonical_bonds}                                           \\
        (\mathcal R\mathcal N_{A^0}(B))(\beta)
                            & =|G|^{-n}\sum_q\prod_e B_e(q)_{\beta_e^+,\beta_e^-},\nonumber \\
        c_p(\mathcal N_{A^0}(B))
                            & =|G|^{-n}\sum_q\prod_e\ell_{e,p_e}(B_e(q)).\nonumber
    \end{align}
    These identities require neither semi-regularity nor restrictions on $B$.
\end{theorem}
\begin{proof}\leanok
    \uses{thm:peps_dependent_projector_expansion,
        thm:peps_dependent_coefficient_extraction}
    Separate the $e$-th factor for the slice identity. Regroup
    (\ref{eq:peps_dependent_projector_expansion}) by $J$, obtaining the
    first two formulas. Apply the product-functional identity to each
    term to obtain the third formula.
\end{proof}

\begin{theorem}[An uncut bond forces representation-matrix support]%
    \label{thm:peps_dependent_uncut_support}
    \lean{TNLean.PEPS.DependentBondNetwork.physicalBondSlice_network_averagingSite_mem_range,
        TNLean.PEPS.DependentBondNetwork.physicalBondSlice_mem_range_of_mem_cutSpace}
    \leanok
    \uses{def:peps_dependent_bond_slice, def:peps_dependent_cut}
    If $B_e=I$, then every $S_{e,\tau}\mathcal N_{A^0}(B)$ lies in
    $\operatorname{im}F_e$. Consequently, if $e\notin C$ and
    $\psi\in\mathcal S_{A^0,C}$, then
    $S_{e,\tau}\psi\in\operatorname{im}F_e$ for every $\tau$, even when
    the cut boundary is arbitrarily correlated.
\end{theorem}
\begin{proof}\leanok
    \uses{thm:peps_dependent_canonical_bond_sum, thm:peps_dependent_cut_basis}
    In a summand of (\ref{eq:peps_dependent_canonical_bonds}), an identity
    insertion gives $B_e(q)=U_e(q_{h(e)}q_{t(e)}^{-1})$.
    Every slice is a scalar multiple of this representation matrix; sums
    retain membership in $\operatorname{im}F_e$.
    Expand a general boundary in fixed configurations. Each resulting
    matrix-unit network still has $B_e=I$ when $e\notin C$.
\end{proof}

\begin{theorem}[Complementary cuts give a coherent bond expansion]%
    \label{thm:peps_dependent_cut_expansion}
    \lean{TNLean.PEPS.DependentBondNetwork.bondRegrouping_mem_product_range_of_mem_cutSpaces,
        TNLean.PEPS.DependentBondNetwork.exists_bondCoefficients_of_mem_cutSpaces,
        TNLean.PEPS.DependentBondNetwork.eq_sum_extracted_bondProducts_of_mem_cutSpaces}
    \leanok
    \uses{def:peps_dependent_cut, def:peps_dependent_bond_coefficient}
    Let $(C_i)_{i\in J}$ be any family of cuts such that every bond is uncut
    in at least one $C_i$. If $\psi\in\mathcal S_{A^0,C_i}$ for every $i$,
    then $\mathcal R\psi\in\operatorname{im}(\bigotimes_eF_e)$.
    In particular there are coefficients $a_p$ with $\psi=\sum_pa_pb_p$.
    If all $U_e$ are semi-regular, this becomes
    \begin{align}
        \psi=\sum_{p:\mathcal E\to G}c_p(\psi)b_p.
        \label{eq:peps_dependent_extracted_expansion}
    \end{align}
\end{theorem}
\begin{proof}\leanok
    \uses{thm:peps_dependent_uncut_support, thm:peps_dependent_product_range,
        thm:peps_dependent_coefficient_extraction}
    For each bond, choose a cut leaving it uncut. The preceding theorem
    places every slice in $\operatorname{im}F_e$. The product-range
    criterion supplies a preimage under $\bigotimes_eF_e$, whose entries
    are the coefficients $a_p$. Apply $c_p$ to the whole expansion to
    identify $a_p=c_p(\psi)$ in the semi-regular case.
\end{proof}

\begin{theorem}[A nonzero coefficient has compatible uncut labels]%
    \label{thm:peps_dependent_uncut_coefficients}
    \lean{TNLean.PEPS.DependentBondNetwork.bondCoefficientExtraction_network_eq_zero_of_incompatible,
        TNLean.PEPS.DependentBondNetwork.bondCoefficientExtraction_eq_zero_of_mem_cutSpace_of_incompatible,
        TNLean.PEPS.DependentBondNetwork.exists_vertexLabels_of_bondCoefficientExtraction_ne_zero}
    \leanok
    \uses{def:peps_dependent_bond_coefficient, def:peps_dependent_cut}
    Assume every $U_e$ is semi-regular. If $B_e=I$ outside $C$ and no
    $q:\mathcal V\to G$ satisfies
    $p_e=q_{h(e)}q_{t(e)}^{-1}$ simultaneously for all $e\notin C$, then
    $c_p(\mathcal N_{A^0}(B))=0$. The same vanishing holds for every
    $\psi\in\mathcal S_{A^0,C}$. Equivalently,
    \begin{align}
        \psi\in\mathcal S_{A^0,C},\quad c_p(\psi)\ne0
        \quad\Longrightarrow\quad
        \exists q:\mathcal V\to G,\quad
        \forall e\notin C,\quad p_e=q_{h(e)}q_{t(e)}^{-1}.
        \label{eq:peps_dependent_uncut_coefficients}
    \end{align}
\end{theorem}
\begin{proof}\leanok
    \uses{thm:peps_dependent_canonical_bond_sum,
        thm:peps_dependent_coefficient_extraction, thm:peps_dependent_cut_basis}
    Fix a vertex assignment $q$ in the extracted sum of
    (\ref{eq:peps_dependent_canonical_bonds}). Incompatibility supplies
    an uncut bond with $p_e\ne q_{h(e)}q_{t(e)}^{-1}$; its trace-dual
    factor is zero. Thus every summand vanishes. For an arbitrary boundary,
    expand into fixed-coordinate columns and use their matrix-unit
    realization. The implication is the contrapositive of this vanishing.
\end{proof}
